\documentclass[12pt, a4paper]{article}

\usepackage{geometry}
\geometry{a4paper, top=2.5cm, bottom=2.5cm, left=2.5cm, right=3cm}
\usepackage{graphicx}
\usepackage{booktabs}
\usepackage{longtable}
\usepackage{array}
\usepackage{multirow}
\usepackage{xcolor}
\usepackage{hyperref}
\usepackage{fancyhdr}
\usepackage{enumitem}
\usepackage{fontspec}
\usepackage{xepersian}
\settextfont[Scale=1.1]{Vazirmatn}
\setlatintextfont{Times New Roman}

\definecolor{primaryorange}{RGB}{234, 88, 12}
\definecolor{lightgray}{RGB}{245, 245, 245}
\definecolor{darkgray}{RGB}{80, 80, 80}

\hypersetup{
    colorlinks=true,
    linkcolor=primaryorange,
    urlcolor=primaryorange
}

\pagestyle{fancy}
\fancyhf{}
\rhead{پروژه پاتوق — سامانه رزرو رستوران}
\lhead{گزارش \lr{User Story}ها}
\rfoot{\thepage}
\renewcommand{\headrulewidth}{0.4pt}

\begin{document}

% \pagenumbering{persian}  % حذف شد برای سازگاری با hyperref

%% ─── صفحه عنوان ────────────────────────────────────────────────────────────
\begin{titlepage}
\centering
\vspace*{2cm}
{\LARGE\bfseries پروژه پاتوق (پاتوق)\\[0.5em]}
{\Large\bfseries سامانه مدیریت و رزرو کافه و رستوران\\[2em]}
\rule{\linewidth}{0.5mm}\\[1em]
{\huge\bfseries\color{primaryorange} گزارش \lr{User Story}های پروژه}\\[1em]
\rule{\linewidth}{0.5mm}\\[2em]
{\large درس: طراحی نرم‌افزار\\[0.5em]}
{\large تاریخ تهیه: تیر ۱۴۰۵}\\[3em]
\vfill
{\normalsize\color{darkgray}
این گزارش \lr{User Story}های تمامی نقش‌های کاربری سیستم پاتوق را
همراه با معیارهای پذیرش، اولویت‌بندی و وضعیت پیاده‌سازی
شرح می‌دهد.}
\end{titlepage}

\tableofcontents
\newpage

%% ─── ۱. مقدمه ──────────────────────────────────────────────────────────────
\section{مقدمه}

\subsection{معرفی سیستم}
پاتوق (\lr{Patogh}) یک سامانه جامع مدیریت و رزرو رستوران است که به کاربران امکان می‌دهد
رستوران‌های مورد علاقه خود را جستجو کرده، میز رزرو کنند و تجربه غذاخوری خود را مدیریت نمایند.
این سیستم برای درس طراحی نرم‌افزار دانشگاهی توسعه یافته است.

\subsection{معماری کلی}
\begin{itemize}[rightmargin=0.5cm]
  \item \textbf{بک‌اند:} \lr{ASP.NET Core 8} با معماری تمیز (\lr{Clean Architecture})
  \item \textbf{فرانت‌اند:} \lr{React 18 + TypeScript + Vite}
  \item \textbf{پایگاه داده:} \lr{PostgreSQL} از طریق \lr{EF Core} + \lr{Redis} برای \lr{OTP} و کش
  \item \textbf{احراز هویت:} \lr{JWT} دوگانه (توکن دسترسی + توکن تازه‌سازی) و ورود با \lr{OTP}
\end{itemize}

\subsection{نقش‌های کاربری}
سیستم سه نقش اصلی دارد:
\begin{enumerate}[rightmargin=0.5cm]
  \item \textbf{مشتری (\lr{Customer}):} جستجو، مشاهده رستوران، رزرو میز، مدیریت رزروها
  \item \textbf{مدیر رستوران (\lr{RestaurantOwner}):} ثبت و مدیریت رستوران، منو، میز و رزروها
  \item \textbf{مدیر ارشد (\lr{Admin}):} تأیید رستوران‌ها، مدیریت کاربران، نظارت کلان
\end{enumerate}

\newpage

%% ─── ۲. User Story‌های مشتری ───────────────────────────────────────────────
\section{\lr{User Story}های نقش مشتری}

\subsection{US-C01: ورود با کد یکبار مصرف (OTP)}
\begin{description}[rightmargin=0.5cm]
  \item[عنوان:] ورود به سیستم با شماره موبایل و کد \lr{OTP}
  \item[داستان:] به عنوان مشتری، می‌خواهم با وارد کردن شماره موبایلم و دریافت کد \lr{OTP}
  وارد سیستم شوم تا نیازی به یادآوری رمز عبور نداشته باشم.
  \item[معیارهای پذیرش:]
  \begin{itemize}[rightmargin=0.5cm]
    \item کد \lr{OTP} ظرف ۳۰ ثانیه دریافت می‌شود
    \item کد دارای اعتبار ۲ دقیقه‌ای است
    \item کد پس از استفاده یکبار باطل می‌شود
    \item در صورت ورود اشتباه، پیام خطای مناسب فارسی نمایش داده می‌شود
    \item سیستم پس از ۵ درخواست در دقیقه، دسترسی را محدود می‌کند
  \end{itemize}
  \item[اولویت:] بسیار بالا \quad \textbf{وضعیت:} پیاده‌سازی شده
\end{description}

\subsection{US-C02: ثبت‌نام در سیستم}
\begin{description}[rightmargin=0.5cm]
  \item[عنوان:] ثبت‌نام با شماره موبایل و انتخاب نقش
  \item[داستان:] به عنوان کاربر جدید، می‌خواهم با شماره موبایل ثبت‌نام کنم و نقش مشتری
  یا مدیر رستوران را انتخاب کنم.
  \item[معیارهای پذیرش:]
  \begin{itemize}[rightmargin=0.5cm]
    \item صفحه انتخاب نقش قبل از ورود نمایش داده می‌شود
    \item اعتبارسنجی فرمت شماره موبایل ایرانی
    \item پیام خوش‌آمدگویی پس از ثبت‌نام موفق
    \item هدایت خودکار به صفحه تکمیل پروفایل برای کاربران جدید
  \end{itemize}
  \item[اولویت:] بسیار بالا \quad \textbf{وضعیت:} پیاده‌سازی شده
\end{description}

\subsection{US-C03: جستجو و کشف رستوران}
\begin{description}[rightmargin=0.5cm]
  \item[عنوان:] جستجوی رستوران بر اساس نام و نوع غذا
  \item[داستان:] به عنوان مشتری، می‌خواهم رستوران‌ها را بر اساس نام، نوع غذا یا موقعیت
  جغرافیایی جستجو کنم.
  \item[معیارهای پذیرش:]
  \begin{itemize}[rightmargin=0.5cm]
    \item نمایش کارت رستوران با عکس، نام، نوع غذا و امتیاز
    \item فیلتر کردن بر اساس نوع غذا
    \item نمایش فقط رستوران‌های تأییدشده
    \item صفحه‌بندی نتایج جستجو
  \end{itemize}
  \item[اولویت:] بالا \quad \textbf{وضعیت:} پیاده‌سازی شده
\end{description}

\subsection{US-C04: مشاهده جزئیات رستوران}
\begin{description}[rightmargin=0.5cm]
  \item[عنوان:] مشاهده اطلاعات کامل رستوران
  \item[داستان:] به عنوان مشتری، می‌خواهم صفحه کامل رستوران را با منو، تصاویر، ساعات
  کاری و میزهای موجود ببینم.
  \item[معیارهای پذیرش:]
  \begin{itemize}[rightmargin=0.5cm]
    \item نمایش منوی کامل با قیمت‌ها
    \item نمایش ساعات کاری هفتگی
    \item دکمه رزرو میز فعال و قابل کلیک
    \item نمایش تصاویر رستوران
  \end{itemize}
  \item[اولویت:] بالا \quad \textbf{وضعیت:} پیاده‌سازی شده
\end{description}

\subsection{US-C05: رزرو میز}
\begin{description}[rightmargin=0.5cm]
  \item[عنوان:] رزرو میز در رستوران
  \item[داستان:] به عنوان مشتری، می‌خواهم میز مناسب برای تعداد نفرات و تاریخ مورد نظرم
  را رزرو کنم.
  \item[معیارهای پذیرش:]
  \begin{itemize}[rightmargin=0.5cm]
    \item انتخاب تاریخ (با تقویم شمسی) و ساعت شروع
    \item نمایش میزهای در دسترس برای آن بازه زمانی
    \item مدت رزرو ثابت ۲ ساعته (قانون کسب‌وکار)
    \item تأیید رزرو با نمایش خلاصه اطلاعات
    \item ارسال اعلان پس از ثبت رزرو
  \end{itemize}
  \item[اولویت:] بسیار بالا \quad \textbf{وضعیت:} پیاده‌سازی شده
\end{description}

\subsection{US-C06: مدیریت رزروها}
\begin{description}[rightmargin=0.5cm]
  \item[عنوان:] مشاهده و لغو رزروها
  \item[داستان:] به عنوان مشتری، می‌خواهم لیست رزروهای گذشته و آینده‌ام را ببینم و در
  صورت لزوم آن‌ها را لغو کنم.
  \item[معیارهای پذیرش:]
  \begin{itemize}[rightmargin=0.5cm]
    \item نمایش رزروها با وضعیت‌های مختلف (در انتظار، تأیید شده، لغو شده)
    \item امکان لغو رزرو قبل از زمان مقرر
    \item نمایش پیام تأیید قبل از لغو
    \item به‌روزرسانی فوری وضعیت پس از لغو
  \end{itemize}
  \item[اولویت:] بالا \quad \textbf{وضعیت:} پیاده‌سازی شده
\end{description}

\subsection{US-C07: علاقه‌مندی‌ها}
\begin{description}[rightmargin=0.5cm]
  \item[عنوان:] ذخیره رستوران‌های مورد علاقه
  \item[داستان:] به عنوان مشتری، می‌خواهم رستوران‌های مورد علاقه‌ام را نشانه‌گذاری کرده
  و بعداً به آن‌ها دسترسی سریع داشته باشم.
  \item[معیارهای پذیرش:]
  \begin{itemize}[rightmargin=0.5cm]
    \item دکمه قلب برای افزودن/حذف از علاقه‌مندی‌ها
    \item صفحه مستقل نمایش علاقه‌مندی‌ها
    \item تغییر فوری وضعیت بدون بارگذاری مجدد صفحه
  \end{itemize}
  \item[اولویت:] متوسط \quad \textbf{وضعیت:} پیاده‌سازی شده
\end{description}

\subsection{US-C08: اعلان‌ها}
\begin{description}[rightmargin=0.5cm]
  \item[عنوان:] دریافت اعلان‌های مربوط به رزرو
  \item[داستان:] به عنوان مشتری، می‌خواهم اعلان‌های مرتبط با وضعیت رزروهایم را دریافت کنم.
  \item[معیارهای پذیرش:]
  \begin{itemize}[rightmargin=0.5cm]
    \item اعلان تأیید رزرو توسط رستوران
    \item اعلان رد رزرو با دلیل
    \item اعلان یادآوری قبل از وقت رزرو
    \item نمایش تعداد اعلان‌های خوانده‌نشده
  \end{itemize}
  \item[اولویت:] متوسط \quad \textbf{وضعیت:} پیاده‌سازی شده
\end{description}

\subsection{US-C09: مدیریت پروفایل}
\begin{description}[rightmargin=0.5cm]
  \item[عنوان:] ویرایش اطلاعات شخصی و آپلود آواتار
  \item[داستان:] به عنوان مشتری، می‌خواهم نام نمایشی و تصویر پروفایل خود را تنظیم کنم.
  \item[معیارهای پذیرش:]
  \begin{itemize}[rightmargin=0.5cm]
    \item ویرایش نام نمایشی
    \item آپلود و تغییر تصویر پروفایل
    \item تغییر رمز عبور
    \item تأیید تغییرات با پیام موفقیت
  \end{itemize}
  \item[اولویت:] پایین \quad \textbf{وضعیت:} پیاده‌سازی شده
\end{description}

\newpage

%% ─── ۳. User Story‌های مدیر رستوران ─────────────────────────────────────────
\section{\lr{User Story}های نقش مدیر رستوران}

\subsection{US-M01: ثبت رستوران جدید}
\begin{description}[rightmargin=0.5cm]
  \item[عنوان:] ثبت رستوران با فرایند چندمرحله‌ای
  \item[داستان:] به عنوان مدیر رستوران، می‌خواهم اطلاعات کامل رستورانم را در چندین مرحله
  وارد کنم تا برای تأیید ادمین ارسال شود.
  \item[معیارهای پذیرش:]
  \begin{itemize}[rightmargin=0.5cm]
    \item فرم چندمرحله‌ای (نام، توضیح، موقعیت، نوع غذا، ساعات کاری، تصاویر)
    \item آپلود تصاویر رستوران
    \item ذخیره پیش‌نویس در هر مرحله
    \item وضعیت اولیه: در انتظار تأیید (\lr{Pending})
  \end{itemize}
  \item[اولویت:] بسیار بالا \quad \textbf{وضعیت:} پیاده‌سازی شده
\end{description}

\subsection{US-M02: مدیریت منوی رستوران}
\begin{description}[rightmargin=0.5cm]
  \item[عنوان:] افزودن، ویرایش و حذف آیتم‌های منو
  \item[داستان:] به عنوان مدیر رستوران، می‌خواهم منوی رستورانم را به‌روز نگه دارم.
  \item[معیارهای پذیرش:]
  \begin{itemize}[rightmargin=0.5cm]
    \item افزودن آیتم جدید با نام، توضیح، قیمت و تصویر
    \item ویرایش آیتم‌های موجود
    \item حذف آیتم با تأیید
    \item نمایش لیست کامل منو
  \end{itemize}
  \item[اولویت:] بالا \quad \textbf{وضعیت:} پیاده‌سازی شده
\end{description}

\subsection{US-M03: مدیریت میزها}
\begin{description}[rightmargin=0.5cm]
  \item[عنوان:] تعریف و مدیریت میزهای رستوران
  \item[داستان:] به عنوان مدیر رستوران، می‌خواهم تعداد و ظرفیت میزهای رستورانم را
  تعریف و مدیریت کنم.
  \item[معیارهای پذیرش:]
  \begin{itemize}[rightmargin=0.5cm]
    \item افزودن میز با شماره و ظرفیت
    \item حذف میز
    \item نمایش وضعیت اشغال بودن میز
  \end{itemize}
  \item[اولویت:] بالا \quad \textbf{وضعیت:} پیاده‌سازی شده
\end{description}

\subsection{US-M04: مدیریت رزروهای رستوران}
\begin{description}[rightmargin=0.5cm]
  \item[عنوان:] مشاهده، تأیید و رد رزروها
  \item[داستان:] به عنوان مدیر رستوران، می‌خواهم رزروهای ثبت‌شده را ببینم و آن‌ها را
  تأیید یا رد کنم.
  \item[معیارهای پذیرش:]
  \begin{itemize}[rightmargin=0.5cm]
    \item نمایش لیست رزروها با فیلتر بر اساس تاریخ و وضعیت
    \item تأیید رزرو با ارسال اعلان به مشتری
    \item رد رزرو با امکان وارد کردن دلیل
    \item لغو رزرو توسط مدیر رستوران
  \end{itemize}
  \item[اولویت:] بسیار بالا \quad \textbf{وضعیت:} پیاده‌سازی شده
\end{description}

\subsection{US-M05: تنظیم ساعات کاری}
\begin{description}[rightmargin=0.5cm]
  \item[عنوان:] تنظیم ساعات کاری هفتگی رستوران
  \item[داستان:] به عنوان مدیر رستوران، می‌خواهم ساعات باز و بسته بودن رستوران را برای
  هر روز هفته به‌طور جداگانه تنظیم کنم.
  \item[معیارهای پذیرش:]
  \begin{itemize}[rightmargin=0.5cm]
    \item تنظیم ساعت شروع و پایان برای هر روز هفته
    \item علامت‌گذاری روزهای تعطیل
    \item ذخیره تغییرات با تأیید
  \end{itemize}
  \item[اولویت:] متوسط \quad \textbf{وضعیت:} پیاده‌سازی شده
\end{description}

\subsection{US-M06: داشبورد مدیر رستوران}
\begin{description}[rightmargin=0.5cm]
  \item[عنوان:] مشاهده آمار و وضعیت کلی رستوران
  \item[داستان:] به عنوان مدیر رستوران، می‌خواهم در یک نگاه آمار رزروها و وضعیت امروز
  رستورانم را ببینم.
  \item[معیارهای پذیرش:]
  \begin{itemize}[rightmargin=0.5cm]
    \item نمایش تعداد کل رزروها، تأییدشده‌ها و در انتظار
    \item نمایش رزروهای امروز
    \item نمودار وضعیت رزروها
  \end{itemize}
  \item[اولویت:] متوسط \quad \textbf{وضعیت:} پیاده‌سازی شده
\end{description}

\newpage

%% ─── ۴. User Story‌های مدیر ارشد ────────────────────────────────────────────
\section{\lr{User Story}های نقش مدیر ارشد (Admin)}

\subsection{US-A01: ورود امن مدیر ارشد}
\begin{description}[rightmargin=0.5cm]
  \item[عنوان:] ورود مدیر ارشد با نام کاربری و رمز عبور
  \item[داستان:] به عنوان مدیر ارشد، می‌خواهم از طریق صفحه ورود مخصوص ادمین با اعتبارنامه
  امن وارد سیستم شوم.
  \item[معیارهای پذیرش:]
  \begin{itemize}[rightmargin=0.5cm]
    \item صفحه ورود جداگانه در \lr{/admin/login}
    \item احراز هویت با رمز عبور هش‌شده (بدون \lr{OTP})
    \item محافظت در برابر حملات \lr{Brute Force}
    \item هدایت خودکار به داشبورد پس از ورود موفق
  \end{itemize}
  \item[اولویت:] بسیار بالا \quad \textbf{وضعیت:} پیاده‌سازی شده
\end{description}

\subsection{US-A02: تأیید رستوران‌ها}
\begin{description}[rightmargin=0.5cm]
  \item[عنوان:] بررسی و تأیید یا رد درخواست‌های ثبت رستوران
  \item[داستان:] به عنوان مدیر ارشد، می‌خواهم درخواست‌های ثبت رستوران جدید را بررسی
  کرده و آن‌ها را تأیید یا رد کنم.
  \item[معیارهای پذیرش:]
  \begin{itemize}[rightmargin=0.5cm]
    \item نمایش لیست رستوران‌های در انتظار تأیید
    \item مشاهده جزئیات کامل هر رستوران
    \item تأیید با یک کلیک و ارسال اعلان به مدیر رستوران
    \item رد با امکان وارد کردن دلیل
  \end{itemize}
  \item[اولویت:] بسیار بالا \quad \textbf{وضعیت:} پیاده‌سازی شده
\end{description}

\subsection{US-A03: مدیریت کاربران}
\begin{description}[rightmargin=0.5cm]
  \item[عنوان:] مشاهده لیست کاربران و تغییر نقش
  \item[داستان:] به عنوان مدیر ارشد، می‌خواهم لیست تمام کاربران را ببینم و در صورت لزوم
  نقش آن‌ها را تغییر دهم.
  \item[معیارهای پذیرش:]
  \begin{itemize}[rightmargin=0.5cm]
    \item نمایش لیست کاربران با صفحه‌بندی
    \item فیلتر بر اساس نقش
    \item تغییر نقش کاربر
    \item نمایش اطلاعات تماس کاربر
  \end{itemize}
  \item[اولویت:] متوسط \quad \textbf{وضعیت:} پیاده‌سازی شده
\end{description}

\subsection{US-A04: داشبورد آمار کلی}
\begin{description}[rightmargin=0.5cm]
  \item[عنوان:] مشاهده آمار کلی سیستم
  \item[داستان:] به عنوان مدیر ارشد، می‌خواهم در داشبورد خود آمار کلی سیستم شامل
  تعداد کاربران، رستوران‌ها و رزروها را ببینم.
  \item[معیارهای پذیرش:]
  \begin{itemize}[rightmargin=0.5cm]
    \item تعداد کل کاربران، رستوران‌ها و رزروها
    \item تعداد رستوران‌های در انتظار تأیید
    \item نمودار رشد ماهانه
  \end{itemize}
  \item[اولویت:] متوسط \quad \textbf{وضعیت:} پیاده‌سازی شده
\end{description}

\newpage

%% ─── ۵. جدول اولویت‌بندی ────────────────────────────────────────────────────
\section{جدول اولویت‌بندی \lr{User Story}ها}

\begin{longtable}{|r|r|r|r|r|}
\hline
\textbf{کد} & \textbf{عنوان} & \textbf{نقش} & \textbf{اولویت} & \textbf{وضعیت} \\
\hline
\endhead
\lr{US-C01} & ورود با \lr{OTP} & مشتری & بسیار بالا & انجام شده \\
\hline
\lr{US-C02} & ثبت‌نام & مشتری & بسیار بالا & انجام شده \\
\hline
\lr{US-C05} & رزرو میز & مشتری & بسیار بالا & انجام شده \\
\hline
\lr{US-M01} & ثبت رستوران & مدیر رستوران & بسیار بالا & انجام شده \\
\hline
\lr{US-M04} & مدیریت رزروها & مدیر رستوران & بسیار بالا & انجام شده \\
\hline
\lr{US-A01} & ورود ادمین & مدیر ارشد & بسیار بالا & انجام شده \\
\hline
\lr{US-A02} & تأیید رستوران & مدیر ارشد & بسیار بالا & انجام شده \\
\hline
\lr{US-C03} & جستجوی رستوران & مشتری & بالا & انجام شده \\
\hline
\lr{US-C04} & جزئیات رستوران & مشتری & بالا & انجام شده \\
\hline
\lr{US-C06} & مدیریت رزروها & مشتری & بالا & انجام شده \\
\hline
\lr{US-M02} & مدیریت منو & مدیر رستوران & بالا & انجام شده \\
\hline
\lr{US-M03} & مدیریت میزها & مدیر رستوران & بالا & انجام شده \\
\hline
\lr{US-C07} & علاقه‌مندی‌ها & مشتری & متوسط & انجام شده \\
\hline
\lr{US-C08} & اعلان‌ها & مشتری & متوسط & انجام شده \\
\hline
\lr{US-M05} & ساعات کاری & مدیر رستوران & متوسط & انجام شده \\
\hline
\lr{US-M06} & داشبورد مدیر & مدیر رستوران & متوسط & انجام شده \\
\hline
\lr{US-A03} & مدیریت کاربران & مدیر ارشد & متوسط & انجام شده \\
\hline
\lr{US-A04} & آمار کلی & مدیر ارشد & متوسط & انجام شده \\
\hline
\lr{US-C09} & مدیریت پروفایل & مشتری & پایین & انجام شده \\
\hline
\end{longtable}

\newpage

%% ─── ۶. وابستگی‌ها ──────────────────────────────────────────────────────────
\section{وابستگی‌های بین \lr{User Story}ها}

\begin{itemize}[rightmargin=0.5cm]
  \item \lr{US-C05} (رزرو میز) به \lr{US-C01} (ورود) و \lr{US-C04} (جزئیات رستوران) وابسته است.
  \item \lr{US-C06} (مدیریت رزروها) به \lr{US-C05} (رزرو میز) وابسته است.
  \item \lr{US-C07} (علاقه‌مندی‌ها) به \lr{US-C01} (ورود) وابسته است.
  \item \lr{US-M04} (مدیریت رزروها مدیر) به \lr{US-M01} (ثبت رستوران) وابسته است.
  \item \lr{US-M02} و \lr{US-M03} به \lr{US-M01} وابسته هستند.
  \item \lr{US-A02} (تأیید رستوران) پیش‌نیاز نمایش رستوران در لیست عمومی است.
  \item \lr{US-C03} و \lr{US-C04} تنها رستوران‌های تأییدشده توسط \lr{US-A02} را نمایش می‌دهند.
\end{itemize}

\section{نتیجه‌گیری}

تمامی ۱۹ \lr{User Story} تعریف‌شده در این گزارش در طول سه اسپرینت پیاده‌سازی و تحویل داده شدند.
سیستم پاتوق به عنوان یک نمونه کامل از سامانه رزرو رستوران، تمامی قابلیت‌های اساسی را پوشش می‌دهد:
از ثبت‌نام کاربر تا رزرو میز، از مدیریت منو تا تأیید رستوران توسط ادمین.

\end{document}
